\documentclass[border=10pt]{standalone}
\usepackage{fontspec}
\AtBeginDocument{\hyphenation{zenodo}}% never hyphenate "zenodo" inside a DOI
\setmainfont{Noto Serif}[Ligatures=TeXOff]
\XeTeXgenerateactualtext=1
\newfontfamily\bengalifont{Noto Serif Bengali}[Script=Bengali, NFSSFamily=bnfam]
\usepackage[Latin,Bengali]{ucharclasses}
% group-free font switching (as in header-common.tex): robust inside braces, cells and quotes
\makeatletter
\def\BN@saved{\rmdefault}
\def\BN@bn{bnfam}
\newcommand\BNin{\edef\BN@tmp{\f@family}\ifx\BN@tmp\BN@bn\else\let\BN@saved\BN@tmp\fi\fontfamily{bnfam}\selectfont}
\newcommand\BNout{\fontfamily{\BN@saved}\selectfont}
\makeatother
\setTransitionsFor{Bengali}{\BNin}{\BNout}
% Typographic quotes and dashes are not in a ucharclasses class (class 0), so a Bangla word right
% after “ or ‘ would not trigger the Bangla font. Treat them as Latin so that every switch fires.
\XeTeXcharclass"2018=\XeTeXcharclass`A
\XeTeXcharclass"2019=\XeTeXcharclass`A
\XeTeXcharclass"201C=\XeTeXcharclass`A
\XeTeXcharclass"201D=\XeTeXcharclass`A
\XeTeXcharclass"2013=\XeTeXcharclass`A
\XeTeXcharclass"2014=\XeTeXcharclass`A
\XeTeXcharclass"2026=\XeTeXcharclass`A
\usepackage{newunicodechar}
\newunicodechar{।}{{\bengalifont\char"0964}}
\newunicodechar{॥}{{\bengalifont\char"0965}}

\usepackage[table]{xcolor}
\usepackage{booktabs,tabularx,array}
\definecolor{bapilrule}{HTML}{1F2937}
\definecolor{bapilstripe}{HTML}{F3F4F6}
\renewcommand{\arraystretch}{1.28}
\setlength{\tabcolsep}{5pt}

\begin{document}
\begin{minipage}{22cm}
{\large\bfseries সারণি ৭০।}\ {\large শর্তসাপেক্ষ জাতীয় সম্প্রসারণের নির্দেশক বাজেট, বছর ৪ থেকে ৮ (কোটি টাকা)}\par\vspace{7pt}
\small
\rowcolors{2}{bapilstripe}{white}
\begin{tabularx}{\linewidth}{@{}>{\hsize=2.674\hsize\raggedright\arraybackslash}X>{\hsize=0.708\hsize\raggedright\arraybackslash}X>{\hsize=0.708\hsize\raggedright\arraybackslash}X>{\hsize=0.708\hsize\raggedright\arraybackslash}X>{\hsize=0.708\hsize\raggedright\arraybackslash}X>{\hsize=0.708\hsize\raggedright\arraybackslash}X>{\hsize=0.787\hsize\raggedright\arraybackslash}X@{}}
\arrayrulecolor{bapilrule}\toprule
\rowcolor{white}\textbf{বিষয়} & \textbf{বছর ৪} & \textbf{বছর ৫} & \textbf{বছর ৬} & \textbf{বছর ৭} & \textbf{বছর ৮} & \textbf{মোট} \\
\midrule
কেন্দ্রীয় ল্যাবের জনবল & ৫.৭১ & ৭.৪৮ & ১০.০৯ & ১২.৩৮ & ১৪.০৯ & ৪৯.৭৪ \\
বিভাগীয় বিশ্ববিদ্যালয়-কেন্দ্র & ২.৫০ & ৫.৪১ & ৫.৮৪ & ৬.৩১ & ৬.৮১ & ২৬.৮৮ \\
সংস্থায় নিযুক্ত বিশ্লেষক & ১.৩৬ & ২.৪৫ & ৩.৭০ & ৫.১৪ & ৬.১৭ & ১৮.৮২ \\
বাংলা ভাষা-ডেটা ও টীকাকরণ কর্মী & ১.৪০ & ২.৪৫ & ৩.০৪ & ৩.২৮ & ৩.৫৫ & ১৩.৭২ \\
স্নাতক ফেলোশিপ & ১.২১ & ১.৯৬ & ২.৮২ & ৩.০৫ & ৩.২৯ & ১২.৩৩ \\
কম্পিউটিং, হোস্টিং ও সমন্বিত অনুসন্ধান & ৬.২৬ & ৬.৫১ & ৭.০৮ & ৭.২০ & ৯.৫৬ & ৩৬.৬১ \\
ডেটা ও বেঞ্চমার্ক & ১.১০ & ১.৩৬ & ১.৪৭ & ১.৪৩ & ১.৫৪ & ৬.৯০ \\
কর্মকর্তাদের প্রশিক্ষণ & ১.২২ & ১.৯৩ & ২.৫৩ & ৩.২১ & ৪.২৩ & ১৩.১২ \\
শাসন, নিরীক্ষা ও মূল্যায়ন & ০.৪৮ & ০.৫২ & ১.১৫ & ০.৬০ & ১.৫১ & ৪.২৫ \\
স্বাগতিকের পরোক্ষ ব্যয় (১০\%) & ২.১২ & ৩.০১ & ৩.৭৭ & ৪.২৬ & ৫.০৮ & ১৮.২৪ \\
আপৎকালীন সংস্থান (১০\%) & ২.১২ & ৩.০১ & ৩.৭৭ & ৪.২৬ & ৫.০৮ & ১৮.২৪ \\
\textbf{মোট, বছর ৪ থেকে ৮} & \textbf{২৫.৫০} & \textbf{৩৬.০৮} & \textbf{৪৫.২৬} & \textbf{৫১.১১} & \textbf{৬০.৯০} & \textbf{২১৮.৮৫} \\
\bottomrule
\end{tabularx}
\par\vspace{5pt}{\footnotesize\itshape ৩৬ মাসের পরিকল্পনার জন্য অনুরোধকৃত অর্থের অংশ নয়। সারণি ৪৯-এর একই একক-ব্যয়, প্রথম বছরের মূল্যে, বছরে ৮\% মূল্যবৃদ্ধিসহ; ১০\% স্বাগতিকের পরোক্ষ ব্যয় ও ১০\% আপৎকালীন সংস্থান। কেন্দ্রীয় ল্যাব ২৫ থেকে ৪৫টি পদে বাড়ে; প্রতিটি পাঁচটি পদসহ চারটি বিভাগীয় বিশ্ববিদ্যালয়-কেন্দ্র; গ্রহণকারী সংস্থায় ৬ থেকে ২০ জন বিশ্লেষক; ১৭ থেকে ৩৩ জন বাংলা ভাষা-ডেটা কর্মী (টীকাকারের মাসিক ব্যয় ৪৫,০০০ টাকা); বছরে ২০ থেকে ৪০ জন স্নাতক ফেলো; ৫,০০০ কর্মকর্তার প্রশিক্ষণ; জাতীয় ইনফারেন্স ও অনুসন্ধান প্ল্যাটফর্ম এবং সমন্বিত অনুসন্ধানে যুক্ত হওয়া সংস্থাগুলোর সহায়তা (প্রতি সংস্থায় ৫০ লাখ টাকা)। blueprint/\allowbreak{}src/\allowbreak{}budget\_model.py (scaleup ফাংশন) দ্বারা হিসাবকৃত; ব্লুপ্রিন্ট ২.০-এ একটি ব্যয়-নির্ধারিত পরিকল্পনা এর স্থলাভিষিক্ত হবে।}
\par\vspace{3pt}{\scriptsize\color{gray} উৎস: BAPIL টেকনিক্যাল ব্লুপ্রিন্ট, সংস্করণ ১.০ (২০২৬)। CC BY 4.0।}
\end{minipage}
\end{document}
